% SCOUT probe structure (b): 2 views + reverse mapping + (implicit) intersection + energy peak
\begin{tikzpicture}[
    x=1cm, y=1cm, font=\footnotesize, >=Stealth,
    lbl/.style    = {font=\scriptsize, align=center, inner sep=1.5pt},
    note/.style   = {font=\scriptsize, text=black!60},
    arrlbl/.style = {font=\scriptsize, align=center, fill=white,
                     fill opacity=.92, text opacity=1, inner sep=1.5pt},
    segbox/.style = {draw, rounded corners=1pt, minimum width=0.34cm,
                     minimum height=0.30cm, inner sep=0pt, font=\tiny},
  ]

  % ----------------------------------------------------------------
  % Geometry knobs
  % ----------------------------------------------------------------
  \def\mvX{0.65}
  \def\mvN{12}
  \def\mvW{0.42}
  \def\mvH{0.30}
  \def\mvGap{0.70}
  \pgfmathsetmacro{\stripRight}{\mvX+\mvN*\mvW}

  % ----------------------------------------------------------------
  % Segment row (top)
  % ----------------------------------------------------------------
  \def\segX{\mvX}
  \def\segY{0.45}
  \def\segN{10}
  \def\segW{0.50}

  % segments: 1..10
  % true hot: 5,6 ; halo (normal): 4,7 ; halo (small): 3,8 ; rest: neutral
  \foreach \i in {0,...,9} {
    \pgfmathtruncatemacro{\sid}{\i+1}
    \pgfmathtruncatemacro{\isHot}{(\sid==5 || \sid==6) ? 1 : 0}
    \pgfmathtruncatemacro{\isHalo}{(\sid==4 || \sid==7) ? 1 : 0}
    \pgfmathtruncatemacro{\isHaloSmall}{(\sid==3 || \sid==8) ? 1 : 0}

    \ifnum\isHot=1
      \node[segbox, fill=orange!60, line width=0.8pt]
        (seg\sid) at ({\segX+\i*\segW},{\segY}) {\sid};
    \else
      \ifnum\isHalo=1
        \node[segbox, fill=orange!20, line width=0.6pt]
          (seg\sid) at ({\segX+\i*\segW},{\segY}) {\sid};
      \else
        \ifnum\isHaloSmall=1
          \node[segbox, fill=orange!10, draw=orange!25, line width=0.5pt]
            (seg\sid) at ({\segX+\i*\segW},{\segY}) {\sid};
        \else
          \node[segbox, fill=black!5, draw=black!25]
            (seg\sid) at ({\segX+\i*\segW},{\segY}) {\sid};
        \fi
      \fi
    \fi
  }

  % ----------------------------------------------------------------
  % Two view strips (middle)
  % ----------------------------------------------------------------
  \def\mvYstart{-0.55}

  % View 1: selected window at sections 2–5
  \pgfmathsetmacro{\yyA}{\mvYstart}
  \draw[thick] (\mvX,\yyA) rectangle (\stripRight,{\yyA+\mvH});
  \foreach \i in {1,...,11} {
    \draw[thin, black!25]
      ({\mvX+\i*\mvW},\yyA) -- ({\mvX+\i*\mvW},{\yyA+\mvH});
  }
  % halo-ish neighbors (dashed)
  \fill[orange!10]
    ({\mvX+0*\mvW},\yyA) rectangle ({\mvX+2*\mvW},{\yyA+\mvH});
  \draw[thick, orange!45!black, dashed]
    ({\mvX+0*\mvW},\yyA) rectangle ({\mvX+2*\mvW},{\yyA+\mvH});
  \fill[orange!10]
    ({\mvX+6*\mvW},\yyA) rectangle ({\mvX+7*\mvW},{\yyA+\mvH});
  \draw[thick, orange!45!black, dashed]
    ({\mvX+6*\mvW},\yyA) rectangle ({\mvX+7*\mvW},{\yyA+\mvH});
  % main selected window (solid)
  \fill[orange!35]
    ({\mvX+2*\mvW},\yyA) rectangle ({\mvX+6*\mvW},{\yyA+\mvH});
  \draw[thick, orange!70!black]
    ({\mvX+2*\mvW},\yyA) rectangle ({\mvX+6*\mvW},{\yyA+\mvH});
  \node[note, anchor=east] at ({\mvX-0.06},{\yyA+\mvH/2}) {view 1};

  % View 2: selected window at sections 6–9
  \pgfmathsetmacro{\yyB}{\mvYstart - \mvGap}
  \draw[thick] (\mvX,\yyB) rectangle (\stripRight,{\yyB+\mvH});
  \foreach \i in {1,...,11} {
    \draw[thin, black!25]
      ({\mvX+\i*\mvW},\yyB) -- ({\mvX+\i*\mvW},{\yyB+\mvH});
  }
  % halo-ish neighbors (dashed)
  \fill[orange!10]
    ({\mvX+5*\mvW},\yyB) rectangle ({\mvX+6*\mvW},{\yyB+\mvH});
  \draw[thick, orange!45!black, dashed]
    ({\mvX+5*\mvW},\yyB) rectangle ({\mvX+6*\mvW},{\yyB+\mvH});
  \fill[orange!10]
    ({\mvX+10*\mvW},\yyB) rectangle ({\mvX+10.8*\mvW},{\yyB+\mvH});
  \draw[thick, orange!45!black, dashed]
    ({\mvX+10*\mvW},\yyB) rectangle ({\mvX+10.8*\mvW},{\yyB+\mvH});
  % main selected window (solid)
  \fill[orange!35]
    ({\mvX+6*\mvW},\yyB) rectangle ({\mvX+10*\mvW},{\yyB+\mvH});
  \draw[thick, orange!70!black]
    ({\mvX+6*\mvW},\yyB) rectangle ({\mvX+10*\mvW},{\yyB+\mvH});
  \node[note, anchor=east] at ({\mvX-0.06},{\yyB+\mvH/2}) {view 2};

  % ----------------------------------------------------------------
  % Reverse-mapping arrows (avoid crossing)
  % ----------------------------------------------------------------
  % View1 arrow -> near seg5
  \coordinate (v1c) at ({\mvX+4*\mvW},{\yyA+\mvH});
  \coordinate (s1c) at ({\segX+4*\segW},{\segY-0.18}); % around seg5
  \draw[->, thick, black!55] (v1c) -- (s1c);

  % View2 arrow -> near seg5
  \coordinate (v2c) at ({\mvX+8*\mvW},{\yyB+\mvH});
  \coordinate (s2c) at ({\segX+4*\segW},{\segY-0.18}); % around seg5
  \draw[->, thick, black!55] (v2c) -- (s2c);

  % ----------------------------------------------------------------
  % Intersection marker above segment row: brace over seg5..6 labeled \mathcal{C}
  % (keep the brace, but move \bigcap_v to below seg6)
  \draw[decorate, decoration={brace, amplitude=3.5pt}]
    (seg5.north west) -- (seg6.north east)
    node[midway, above=4pt, lbl] {$\mathcal{C}$};

  % \bigcap_v moved to below seg6
  \node[note, anchor=north] at ($(seg6.south)+(0.25,-0.02)$) {$\bigcap_v$};

  % ----------------------------------------------------------------
  % Energy spectrum (bottom) — simplified labels
  % ----------------------------------------------------------------
  \pgfmathsetmacro{\barY}{\yyB - 1.15}
  \def\barX{\mvX}
  \def\barW{\mvW}
  \def\barScale{0.85}

  % Bars: peak aligned with view-1 selected window (sections 2–5)
  \foreach \idx/\ht in {
    0/0.08, 1/0.10,
    2/0.65, 3/0.80, 4/0.72, 5/0.50,
    6/0.10, 7/0.07, 8/0.09, 9/0.06, 10/0.08, 11/0.07} {
    \pgfmathsetmacro{\bh}{\ht*\barScale}
    \pgfmathtruncatemacro{\isPeak}{(\idx>=2 && \idx<=5) ? 1 : 0}
    \ifnum\isPeak=1
      \fill[orange!45]
        ({\barX+\idx*\barW},\barY)
        rectangle ({\barX+(\idx+1)*\barW},{\barY+\bh});
      \draw[orange!70!black]
        ({\barX+\idx*\barW},\barY)
        rectangle ({\barX+(\idx+1)*\barW},{\barY+\bh});
    \else
      \fill[black!10]
        ({\barX+\idx*\barW},\barY)
        rectangle ({\barX+(\idx+1)*\barW},{\barY+\bh});
      \draw[black!30]
        ({\barX+\idx*\barW},\barY)
        rectangle ({\barX+(\idx+1)*\barW},{\barY+\bh});
    \fi
  }
  \draw[thick] (\barX,\barY) -- ({\barX+12*\barW},\barY);

  % x-label
  \node[note] at ({\barX+6*\barW},{\barY-0.18}) {section $s$};

  % Move E_s to the right of the peak (avoid overlap with left label)
  \node[note, rotate=90, anchor=south]
    at ({\barX+7.35*\barW},{\barY+0.45*\barScale}) {$E_s$};

  % peak brace with only "peak"
  \draw[decorate, decoration={brace, amplitude=3pt}, orange!65!black]
    ({\barX+2*\barW},{\barY+0.80*\barScale+0.08})
    -- ({\barX+6*\barW},{\barY+0.80*\barScale+0.08})
    node[midway, above=2pt, lbl, text=orange!65!black] {peak};

  % Left label: align nicely (same x as strip labels), no overlap
  \node[note, anchor=east, align=left]
    at ({\mvX+0.75},{\barY+0.50})
    {view 1 energy};

\end{tikzpicture}